% ============================================================
%  GPF Plume Navigation — Results
%  Drop-in section for the Results chapter.
% ============================================================

\subsection{Plume Navigation Results}

\subsubsection{Training Dynamics}

Figure~\ref{fig:gpf_training_curve} shows the evaluation success rate
throughout training along with the grow and prune events recorded by the GPF
agent.  The agent begins with a single hidden layer and a random policy
($\varepsilon = 1.0$), requiring roughly $500$ episodes to acquire any
consistent navigating behavior.

Three grow events occurred during training.  The first grow fired at
episode $2{,}000$, at which point the one-layer baseline had plateaued near
$97.5\%$ success rate and the split-window stagnation criterion
($\text{improvement} < 5\%$) was met.  A subsequent OBD pruning pass
retained $47.9\%$ of the weights in the original hidden layer, yielding a
sparse two-layer network.  The immediate post-prune evaluation
(episode $2{,}500$) measured $91.5\%$, reflecting a temporary disruption
as the agent adapted to the newly sparse weight structure.

Recovery was rapid: by episode $3{,}000$ the two-layer sparse network had
reached $\mathbf{98.0\%}$ success rate, the highest value observed throughout
the entire run.  This best-checkpoint was saved and used for all subsequent
evaluation.

A second grow fired at episode $3{,}000$ (simultaneously with the peak
evaluation), expanding the network to three layers.  The associated second
prune pass operated over two layers and retained only $19.6\%$ of the
combined weight mass—substantially more aggressive than the first prune—due
to the near-identity initialization of the newly added layer, whose weights
carry low saliency before training.  Despite this disruption, the
three-layer network recovered to $95.5\%$ at episode $3{,}500$.

A third grow fired at episode $4{,}000$ ($96.5\%$, four layers, the maximum
depth), triggering a third prune over three layers that retained $9.9\%$ of
the total weight mass.  The four-layer ultra-sparse network sustained
$95.5\%$ through episode $4{,}500$ before degrading to $68.0\%$ by the
final evaluation at episode $5{,}000$, likely due to continued
$\varepsilon$-greedy noise in a nearly converged, low-$\varepsilon$ regime
combined with the extremely sparse representation.  The training trajectory
is summarized in Table~\ref{tab:gpf_training_curve}.

\begin{table}[h]
\centering
\caption{Evaluation success rate and network state throughout training.
  G = grow event; P = prune event (keep ratio in parentheses).}
\label{tab:gpf_training_curve}
\begin{tabular}{rrrll}
\toprule
\textbf{Episode} & \textbf{Success (\%)} & \textbf{Layers} & \textbf{Event} \\
\midrule
  500 & 24.5 & 1 & --- \\
 1000 & 95.5 & 1 & --- \\
 1500 & 40.5 & 1 & --- \\
 2000 & 97.5 & 2 & G: layer $1 \to 2$ \\
 2000 & ---  & 2 & P: kept $47.9\%$ \\
 2500 & 91.5 & 2 & --- \\
 3000 & \textbf{98.0} & 3 & G: layer $2 \to 3$\ \ (best checkpoint) \\
 3000 & ---  & 3 & P: kept $19.6\%$ (2 layers) \\
 3500 & 95.5 & 3 & --- \\
 4000 & 96.5 & 4 & G: layer $3 \to 4$ \\
 4000 & ---  & 4 & P: kept $9.9\%$ (3 layers) \\
 4500 & 95.5 & 4 & --- \\
 5000 & 68.0 & 4 & --- \\
\bottomrule
\end{tabular}
\end{table}

\subsubsection{Effect of Episode Length}

A key finding from the ablation sequence is the critical role of the
training episode budget.  An earlier configuration using $50\,\text{s}$
episodes ($600$ steps) achieved a peak success rate of $94.5\%$; extending
the budget to $100\,\text{s}$ ($1{,}200$ steps) raised the peak to
$98.0\%$, a gain of $3.5$ percentage points.  The improvement is
attributable to the elimination of timeout failures on difficult starting
configurations—episodes in which the agent is initialized far from the
source or in a weak-signal region of the plume.  With a $50\,\text{s}$
budget, the agent cannot complete a full surge--cast--surge cycle before
the episode terminates; doubling the budget allows the full behavioral
strategy to unfold.

\subsubsection{Generalization Evaluation}

To assess the generalization of the best checkpoint beyond the training
distribution, we evaluated the saved $98.0\%$ policy over $100$
independent episodes with distinct random seeds, each generating a unique
source location, initial agent position, and wind realization.  The agent
successfully navigated to within $0.5\,\text{m}$ of the source in
\textbf{94 out of 100} trials ($94\%$ success rate).  The six failures
were all timeout cases in which the agent correctly tracked the plume but
exhausted the $100\,\text{s}$ episode budget before reaching the success
radius—consistent with the hypothesis that the remaining failure mode is
timing rather than navigational strategy.

\subsubsection{Summary}

The GPF agent learned a high-performance navigation policy through
structured growth: a single OBD pruning pass on the initial layer produced
a sparse two-layer representation that achieved $98.0\%$ on the training
distribution and $94\%$ on held-out random initializations.  These results
establish that the GPF grow-and-prune mechanism successfully compresses a
capable policy into a sparse architecture, and that the 100-second episode
budget is sufficient for the agent to develop and execute the multi-phase
olfactory search strategies—surge, cast, and loiter—that characterize
successful biological plume navigation.
